\section{Witness Class and Comparison Protocol}\label{sec:witness}

\subsection{Clause-multiset equivalence}

Let $F$ and $G$ be CNF packages with variable sets $V_F$ and $V_G$.
We write $F\CM G$ when there is a bijection
$\rho:V_F\rightarrow V_G$ such that applying $\rho$ to every literal of every
clause in $F$ yields exactly the clause multiset of $G$, including
multiplicity.

\begin{definition}[M2.1 comparison record]
For a finite case $(n,m)$, an M2.1 comparison record consists of a bundled
package $F_B$, a split package $F_S$, a repair-transport map $\pi$, the
SAT/UNSAT status of both packages, a verified or refuted \CM{} check, and, when
authorized, the raw inclusion-minimal repair families
$\mathcal{R}_B$ and $\mathcal{R}_S$.
\end{definition}

Counts alone do not prove \CM.  Equal variable counts, clause counts, and
clause-length histograms are only necessary conditions.  A positive
classification requires construction and re-verification of $\rho$.  A
negative classification can be established by a violated invariant: if the
packages have different numbers of used variables or clauses, no variable
bijection can make their clause multisets equal.

\subsection{Step~0 classification}

Before repair enumeration, each bundled/split pair receives exactly one
classification:
\begin{description}[style=nextline,leftmargin=2.2em]
  \item[\texttt{equiv\_cm\_persists}.]
  Status agrees and a variable-renaming bijection witnesses \CM.
  \item[\texttt{sat\_equiv\_only}.]
  SAT/UNSAT status agrees, but \CM{} is refuted.
  \item[\texttt{status\_diverges}.]
  The packages have different satisfiability statuses.
  \item[\texttt{inconclusive}.]
  Generation, solving, or structural comparison did not establish one of the
  preceding outcomes.
\end{description}

The gate is asymmetric.  A reviewed
\texttt{equiv\_cm\_persists} result can authorize the strong repair comparison
because the M1.5 control is retained.  A \texttt{sat\_equiv\_only} result can
still be scientifically useful, but it does not automatically authorize repair
enumeration and cannot support a \CM-protected non-canonicity claim.

\subsection{Repair-family comparison}

For a package with active lever universe $L$, a set $R\subseteq L$ is a raw
inclusion-minimal repair when removing all clause families named by $R$ yields
SAT and no strict subset of $R$ does.  Let $\mathcal{R}_B$ and
$\mathcal{R}_S$ denote the bundled and split repair families.
The transport map extends pointwise:
\[
  \pi(R)=\bigcup_{\ell\in R}\pi(\ell),
  \qquad
  \pi(\mathcal{R}_B)=\{\pi(R):R\in\mathcal{R}_B\}.
\]
The repair presentation is non-canonical under the chosen lever refinement
when
\[
  \pi(\mathcal{R}_B)\ne\mathcal{R}_S.
\]
This inequality is a statement about raw lever-level presentation.  It does not
say that one encoding is semantically preferable, nor does it define a
canonical grouping policy.  Rather, it shows that the human-facing repair
vocabulary needs an explicit abstraction contract.

\subsection{Two-track design}

The design separates two axes that would otherwise be conflated:
\begin{itemize}
  \item \textbf{Option D, alternative expansion.}  Keep $n=2$, increase $m$,
  preserve the legacy-style two selector series, and introduce no voter-pair
  selector machinery.
  \item \textbf{Option C, voter expansion.}  Increase $n$ to three, select a
  witness pair, and expose the selected left and right decisive roles
  separately.
\end{itemize}
The first track tests whether the base-case witness survives a change in the
alternative dimension while keeping its structural basis.  The second tests
whether that basis itself remains available once witness selection becomes
nontrivial.

\begin{proposition}[Interpretation discipline]
An Option D repair-family divergence observed after a verified \CM{} check can
be attributed to lever-presentation granularity within the compared packages.
An Option C status agreement after \CM{} failure establishes an encoding
boundary, but does not by itself establish repair-presentation
non-canonicity.
\end{proposition}
\begin{proof}
Under \CM, the packages have identical clauses up to a variable renaming, so
additional logical strength cannot account for a difference caused by how those
clauses are grouped into levers.  When \CM{} fails, unmatched clauses remain a
possible cause of any later repair divergence.  Status agreement removes only
the coarser possibility of a SAT/UNSAT disagreement.
\end{proof}
